%% Déclare un article sur feuille A4
\documentclass[a4paper]{article}

%% Langues et encodages des polices d'écriture
\usepackage[frenchb]{babel}
\usepackage[utf8]{inputenc}
\usepackage[T1]{fontenc}

%%  Choisis la page et les marges 
\usepackage[a4paper,top=3cm,left=3cm,right=3cm,headsep=45pt,marginparwidth=1.75cm]{geometry}

%% Importe les packages dont on a besoin dans le rapport
\usepackage{lastpage}
\usepackage{graphicx}
\usepackage{caption}
\usepackage[colorlinks=true, allcolors=blue]{hyperref}

\usepackage{fancyhdr}

\fancyhf{}
\lhead{
    \includegraphics[scale=0.25]{img/logobyes.png}
    \includegraphics[scale=0.17]{img/logo-bouygues-construction.jpg}
}
\rhead{
    \includegraphics[scale=0.1]{img/logo-univ.png}
    \includegraphics[scale=0.0425]{img/logo-iut.png}
    \includegraphics[scale=0.31]{img/logo-info.png}
}
\renewcommand{\footrulewidth}{0.4pt}
\fancyfoot{
    Pinon Olivier - IUT Annecy - Département Informatique - 2016/2017
    \hfill
    \thepage{} / \pageref*{LastPage}
}


\begin{document}

    \pagestyle{fancy}
    \thispagestyle{empty}
    \noindent
    \begin{minipage}{.5\textwidth}
        Bouygues Energies et Services \\
        49 avenue du Lac du Bourget \\
        73375 Le Bourget du Lac, France
    \end{minipage}
    \begin{minipage}{.5\textwidth}
    \begin{flushright}
        Pinon Olivier \\
        IUT d'Annecy \\
        DUT INFO \\
        Année 2016/2017 \\
    \end{flushright}
    \end{minipage}
    
    \vfill 
    \begin{center}
		\Huge{\textbf{Rapport de stage : }} \\
        \vspace{20pt}
        \Large{Développement d'un module d'habilitation éléctrique en réalité virtuelle sur moteur de jeu Unity}
        
	\end{center}
    \vfill 
    
    Palanca Jérôme  \hfill Damas Luc

 	\newpage 
     % Remerciements 
    \huge \textbf{Remerciements} \vspace{10pt} \\

    \normalsize
    Je souhaite tout d'abord remercier Monsieur Jérôme PALANCA, mon tuteur de stage, pour la confiance, l'écoute et la sympathie dont il a fait preuve à mon égard. Il a su m'aider à comprendre les problèmes de nature inconnue qui m'étaient posés et a su accepter mes propositions tout en me proposant des améliorations. Le suivi quotidien qu'il a effectué m'a permis de constamment garder en tête les objectifs et ainsi de me guider pas à pas dans le projet. \vspace{10pt} \\

    J'adresse aussi mes remerciements à Monsieur Renaud PAYERNE, qui est à l'origine du projet peu commun sur lequel j'ai travaillé, et m'a permis d'obtenir le stage. \vspace{10pt} \\
	Merci également à Monsieur Luc DAMAS, pour son suivi en cours de stage qui m'a donné de précieuses indications sur le contenu et les axes à souligner dans mon rapport et ma soutenance. \vspace{10pt} \\

    Finalement, je tiens à remercier l'ensemble du bureau d'études de Bouygues Energies et Services. Merci à cette équipe qui a su m'accueillir malgré nos travaux très différents, et qui m'a fait passé un très bon stage dans la bonne humeur quotidienne du bureau. Ils ont su me donner un aperçu du monde du travail très plaisant, à la fois sérieux dans les heures de travail, et détendu dans les moments de pause. \vspace{10pt} \\

    Grâce à leur confiance, à leurs réponses et à leur avenance, j'ai pu apprendre beaucoup sur le monde professionnel. Travailler à leur côtés a été un véritable plaisir, et une experience très enrichissante. \\
    
    \newpage
    % Table des matières 
    \tableofcontents

    \newpage 
    % Introduction 
    \huge \textbf{Introduction} \vspace{20pt} \\
    \normalsize
    Afin de finaliser mon diplôme de DUT Informatique à l'IUT d'Annecy, j'ai effectué mon stage d'une durée de 3 mois dans une entreprise de mon choix, qui s'est naturellement porté vers le secteur de l'industrie, puisque je souhaite continuer mes études dans une école d'ingénieur industriel. \\

    Par l'intermédiaire de madame Nathalie Gruson, j'ai répondu à une annonce qui m'a particulièrement intéressé puisqu'elle proposait de travailler sur de la réalité virtuelle en utilisant Unity, un moteur de jeu que j'ai l'habitude d'utiliser, et qui est une technologie que j'affectionne particulièrement. \\

    L'entreprise n'a rien à voir avec le secteur de l'informatique pure, et c'est justement cette différence avec les autres offres de stage qui m'a poussé à candidater, afin de découvrir un autre monde que les sites et applications web et mobiles, que l'on pratique beaucoup à l'IUT. \\
    
    Bouygues Energies et Services m'a proposé de travailler sur un nouveau module, qui permettrait d'aider leurs techniciens à se former à la maintenance de cellules dites HT (Haute Tension), en mettant à profit les nouvelles technologies, notamment la réalité virtuelle. \\

    Dans la première partie, je présenterais l'entreprise qui m'a accueilli. Puis, je développerai le besoin et la mission que l'on m'a assigné. J'expliquerais ensuite la phase d'étude et de réalisation du projet, et je finirai par un bilan de l'experience enrichissante qu'a été ce stage. \\

    \newpage 
    % Présentation générale 
    \section{Présentation générale de l'entreprise}

    \subsection{Le groupe Bouygues et ses filiales}

    Bouygues est un groupe industriel diversifié français fondé en 1952 par Francis Bouygues et dirigé par son fils Martin Bouygues depuis 1989. \\

    Le groupe est structuré autour de trois activités : le secteur BTP (Bouygues Construction,
Bouygues Immobilier et Colas), les réseaux de télécommunication (Bouygues Telecom) et les médias (TF1). Chaque sous-groupe comprend des filales. Bouygues Energies et Services, par exemple, est une filiale du sous-groupe Bouygues Construction, lui même appartenant au groupe Bouygues. Il est à noter que Bouygues Energies et Services a été renommé en 2013 ; Il s'agissait auparavant de l'entreprise ETDE, rachetée par Bouygues Construction en 1984. \\

    schémaGroupeBouygues \\
    
    Le groupe Bouygues est un des leaders du marché dans tous les secteurs où il possède des groupes. En 2015, le chiffre d’affaires du groupe Bouygues s’élève à 32,428 milliards d’euros, et celui de Bouygues Energies Services s'élève à 925 568 300 euros. \\

    \subsection{Produits et services}
 
    \subsubsection{L'activité de Bouygues Energies et Services}
    Bouygues Energies et Services est une filiale dont l'activité est découpée sur 4 marchés principaux : \\
    
    \begin{itemize}

    \item Les infrastructures numériques : \\
        
        Bouygues Energies et Services connecte les territoires afin que tous puissent bénéficier des avantages du tout numérique. Grâce à son savoir-faire et son expertise mises au service de la création et la modernisation des infrastructures de télécommunication et du numérique, l'entreprise relie et dynamise des territoires auparavant isolés. \\

    \item Les villes et collectivités : \\

        La population mondiale devenant de plus en plus urbaine, et le mouvement s'accélérant, la ville de demain doit être plus accueillante, économe, connectée et sécurisée, de façon à allier la qualité de vie avec une empreinte écologique faible. Bouygues Energies et Services est un des acteurs de cette transition qui concerne les collectivités. \\

    \item Le tertiaire et l'industrie: \\

        Afin de les aider à accroître leur compétitivité, ou à optimiser leurs facture énergétique, Bouygues Energies et Services propose des solutions globales aux entreprise industrielles, dans le but de permettre d'améliorer les process mis en place. \\

    \item L'énergie et le transport : \\

        L'entreprise travaille en collaboration avec de grands clients et services publics, en proposant de nouvelles solutions techniques permettant de réduire les consommations et d'augmenter la qualité de nos transports et éclairages. 

    \end{itemize}
    \vspace{5pt}
    
    Le rôle des techniciens de Bouygues Energies et Services est de travailler en collaboration avec d’autres bureaux d’étude dont les spécialités sont le génie climatique et l’architecture, afin de concevoir le bâtiment et ses fonctionnalités avant d’envoyer les plans et de le construire sur le terrain. Fréquemment (presque toutes les semaines), les ingénieurs et chefs de projets sont amenés à se déplacer sur le terrain pour faire un point sur l’avancée des travaux, et discuter de l'implémentation des nouvelles solutions techniques envisagées par les bureaux d'étude. \\
 
    \subsubsection{Les grands projets de l'entreprise}

    Comme expliqué précédemment, l'entreprise travaille sur des grands projets du secteur public. Le domaine d'expertise de Bouygues Energies et Services est l'éléctricité, et notamment les problématiques existantes autour du stockage et du transport d'énergie avec des pertes toujours plus faibles. \\

    Dans le bureau d'étude dans lequel j'ai travaillé, les techniciens et ingénieurs sont en train d'étudier et de mettre en place de nouvelles solutions techniques sur des projets réels comme les stations de métro de Lyon (Lignes 1 et 2), l’autoroute A9 vers Montpellier, ou encore la nouvelle université de Marseille. \\

    \subsection{Implantation et locaux}
    \subsubsection{Bouygues Energies et Services en France et dans le monde}

    Le groupe Bouygues est présent dans plus de 100 pays sur les 5 continents. La société Bouygues Energies et Services, quant à elle, est implantée dans 25 pays d’Europe, d’Afrique et d’Amérique du Nord, suivant 7 zones géographiques (France IDF, France Ouest, France Est/Italie, Royaume Uni, Suisse/Allemagne, Canada, Afrique). Comptant près de 200 implantations sur le territoire Français, dont 35 dans notre zone (France Est/Italie), l’entreprise est donc établie assez largement et dans le monde entier. \\

    imgFrance \\
    imgZoneGeographiqueFrance \\

    \subsubsection{Les locaux de l'entreprise}

    Dans le cadre de mon stage, j’ai travaillé dans les locaux du technopôle Savoie Technolac, plus précisément dans l'un des bureaux d’étude. \\

    imgLocaux \\
    imgBureauDEtudes \\
    imgAtelier \\
    imgSalleReunion \\

    \subsection{Le personnel}

    Le groupe Bouygues entier emploie plus de 100 000 collaborateurs dont 44\% à l’international. Bouygues Energies et Services est une entreprise d'environ 3800 collaborateurs. \\

    Au sein des locaux dans lesquels j'ai travaillé dans le cadre de mon stage, nous étions environ 60 personnes (les effectifs varient d'une semaine sur l'autre du fait des différents projets en cours).  \\

    Le service informatique n'est pas interne à l'entreprise. En effet, le groupe Bouygues a décidé de faire de l'IT Outsourcing, et donc de ne pas engager d'administrateur système dans le personnel de l'entreprise. La gestion informatique est donc effectuée par une autre filiale du groupe Bouygues, qui est appellée Structis. \\

    \subsection{L'équipe de travail}

    Durant la totalité de mon stage, j'ai travaillé en collaboration unique avec mon tuteur, Jérôme Palanca. Ce dernier n'ayant que peu de connaissances techniques sur l'utilisation des moteurs de jeu et le développement orienté objet, c'est lui qui se chargeait de la modélisation 3D des éléments, pendant que j'étais en charge de la conception objet et du développement réel de l'application. \\
    
    Le projet était effectué en parallèle du reste de l'activité du bureau d'études. J'ai eu une dizaine de collègues dans mon bureau de type open space, ce qui m'a permis de voir l'atmosphère de travail dans une équipe et de m'y habituer, mais tous mes collègues travaillaient sur des projets qui n'étaient pas en relation directe avec le mien. \\
    
    % Présentation du besoin
    \section{Le besoin de l'entreprise}
    \subsection{Le module d'habilitation éléctrique}

        L'habilitation éléctrique est un examen visant à certifier qu'un technicien est capable d'effectuer des opérations de maintenance sur un appareil fonctionnant sous haute tension, en restant dans les règles de sécurité mises en place. \\

        Cette habilitation est inhérente à chaque entreprise, et doit être effectuée par l'entreprise qui emploie le technicien pour être valide, même si l'idée reste globalement la même partout. \\

        C'est un examen d'une trentaine de minutes, qui est supervisé par un maître de formation qui fait passer chaque candidat à l'habilitation, un par un. De plus, l'examen nécéssite d'être effectué dans des conditions presques réelles avec des machines qui sont à la fois coûteuses et qu'il faut maintenir. \\

        [ photos installation machines pour l'habilitation ] \\

        Cette habilitation peut être préparée par la lecture d'un manuel de préparation à l'habilitation électrique, dont on m'a confié un exemplaire dans le but de mieux appréhender les aspects techniques, et comprendre ce sur quoi j'allais travailler.

    \subsection{Le besoin d'un nouveau module}

        Le projet qui m'a été proposé est parti de Monsieur Renaud Payerne, et consiste à simuler cette formation avec la technologie de réalité virtuelle, non pas dans le but de remplacer cet examen obligatoire, mais plutôt d'en faciliter la formation. \\

        ------- L'idée est qu'en utilisant des technologies modernes et ludiques, les futurs habilités auront une formation plus efficace. De plus, la formation qui devait s'effectuer un par un pourrait être effectuée par un formateur pour plus de personnes qu'auparavant, réduisant ainsi les moyens et le temps que la formation nécéssite. \\

        ---- Un formateur devrait aussi pouvoir montrer aux futurs habilités ce qui se passe dans la machine (en retirant les différents caches des machines) afin de mieux les aider à comprendre la machine. De plus, une simulation permet de montrer les éléments qui entourent le formé, ce qui favorise la compréhension. \\

    \subsection{Les fonctionnalités de l'application}

    Durant la réalisation du projet, aucun listing de fonctionnalités préalable n'a été fait, car cela nécéssitait trop de temps. L'application en elle-même n'est qu'un protoype visant à démontrer l'efficacité et l'intérêt d'un nouveau module d'habilitation éléctrique. Nous pouvons cependant effectuer une analyse rétrospective de l'application, et définir les acteurs principaux du système et les actions qu'ils peuvent effectuer dans la simulation. \\

    UMLFormé \\

    UMLFormateur \\

    UML \\

    \subsection{Les solutions utilisées}

    Le choix des moyens techniques a été réalisé en amont de mon stage par l'équipe qui a lancé le projet au bureau d'étude de Bouygues Energies et Services. \\ 

    Afin de développer l'application, nous avons choisi d'utiliser le moteur de jeu Unity, pour les performances, la rapidité de développement, la polyvalence et la modularité qui est la sienne. De plus, nous avons pu utiliser des licences gratuites, car l'utilisation que nous avons du moteur n'est pas commerciale. L'utilisation d'un autre moteur, comme Unreal, ne nous a pas paru être un bon choix, notamment parce qu'il nécéssite des compétences en C++, qui n'est pas dans les langages que je maîtrise parfaitement. [Logo Unity / Unreal Engine] \\

    Bien que j'eusse préféré utiliser Linux comme environnement de développement, j'ai eu à travailler sur Windows car les ordinateurs fournis par Structis tournent uniquement sur ce système d'exploitation. L'avantage est que durant tout le développement, je n'ai pas eu à me soucier de la compatibilité entre différents systèmes d'exploitations, même si Unity gère cela de façon transparente. [ Logo Windows / Linux ] \\
    
    Le language C\# a été utilisé pour le scripting de l'application elle-même, puisque c'est celui que j'avais personnellement le plus utilisé avec le moteur de jeu, je connaissais donc bien l'API. De plus, il existe de nombreux comparatifs de performance entre JavaScript et C\#, qui montrent qu'une application développée avec l'API que nous avons utilisé sera bien plus performante. [ Logo CSharp / JS / Benchmarks ] \\

    Le choix de l'IDE que j'ai utilisé, Visual Studio Code, a été fait suite à une recommendation de mon tuteur, que j'ai approuvé puisque beaucoup de mes collègues à l'IUT m'avaient déjà proposé d'apprendre à l'utiliser. Cet outil très puissant, extrêmement modulable, et adapté au développement pour Unity a été une des très bonne découvertes que j'ai pu faire durant mon stage, et je l'utilise désormais comme mon IDE pour tous les languages dans mes projets personnels. Nous avons volontairement décidé de ne pas utiliser Visual Studio pour Windows ; Ce logiciel aurait coûté une licence en plus pour mon projet car Bouygues Energies et Services fait plus d'un million d'euros de chiffre d'affaires, et que j'ai personnellement préféré apprendre à utiliser une technologie que je n'avais pas encore utilisé. [ Logo VSCode / Visual Studio tout court ].\\

    Pour la modélisation 3D, nous avons choisi d'utiliser SketchUp, car c'est le logiciel que mon tuteur savait utilisé, et qu'il est très facile à prendre à main, j'ai donc également pu m'essayer au dessin assisté par ordinateur, ce qui a été un franc succès. [ Logo SketchUp ] \\

    Réalisation du BAPI / Comparatif en vrai \\

    % Etude et réalisation
    \section{Etude et réalisation de l'application}

    \subsection{Planification des tâches}
    
     Dans le cadre du projet, nous avons choisi un cycle de développement suivant une méthode agile, très comparable à la méthode Scrum. Pour cette raison, il n'a pas été possible de faire un diagramme de Gantt pour la planification des tâches en début de projet. Voici alors le diagramme de Gantt (retrospectif) de la planification des tâches du projet :

     [ insérer GanttProjet ]   

    \subsection{Conception du projet}

    La plus grosse difficulté du projet, à mon goût, a été de suivre le cycle de développement agile que nous avons mis en place, sans avoir constitué un réel cahier des charges au début du stage. Nous avons mis en place une trame Trello, mais cette dernière ne permettait pas de définir une conception objet valable. En effet, même si mon tuteur m'a souvent dit de \"voir loin\" pour le développement, pour chaque partie du système, il était possible que le besoin, que je ne pouvais pas prévoir, change d'une heure sur l'autre. La phase de conception a donc été très compliquée pour moi et a entraîné du retard sur la phase de réalisation. \\
    
    Chaque fois qu'un nouveau problème m'était posé, je prenais du temps pour effectuer un croquis, sorte de schéma semblable à de l'UML. Une fois cette première réflexion effectuée, je faisais valider le principe par mon tuteur afin qu'il comprenne comment j'allais résoudre le problème. Je procédais ensuite à la réalisation d'un diagramme UML (très souvent un diagramme de classe). Pour la réalisation des cellules et transformateurs haute tension, qui ont un fonctionnement bien précis, il a été fréquent que je réalise un diagramme d'état. \\ 

    \subsection{Réalisation du projet}

    Beaucoup de retard a été pris sur la réalisation du projet. En effet, la conception n'étant pas possible à effectuer correctement m'a poussé à revenir plusieurs fois sur les mêmes décisions. \\

    La réalisation du projet a été faite suivant une méthode semblable à la méthode Scrum, vue en cours. Chaque jour, je prenais 5 minutes pour discuter avec mon tuteur du travail restant à effectuer, des machines et des interactions qu'il fallait reproduire, et de ce que je pensais réussir à finir dans la journée. Cette sorte de mélée était importante pour garder en tête les objectifs du projet, et me fixer chaque jour un nouveau point à atteindre. \\

    Chaque fois que j'avais un doute, ou besoin d'une information concernant une machine que je ne comprenais pas, ou une interaction à animer, je pouvais contacter mon tuteur de stage par l'intermédiaire de Lync, un système de dialogue intra-entreprise. \\

    Chaque partie du projet terminée était soumise à validation par mon tuteur. Il a été décidé également de mettre en place un système de sauvegarde hebdomadaire sur le serveur interne de l'entreprise, permettant de mettre en place un système de versionning puisque nous n'avons pas utilisé Git à cause du manque de serveur. \\

    Vers la moitié du stage, j'ai eu également une démonstration de mon projet à faire auprès de monsieur Payerne, qui a apprécié le rendu graphique qu'il a vu, et a pu me faire quelques retours. \\

    \subsection{Les tests du projets}

    Unity ne propose pas nativement de système de tests pour les applications qu'il produit. Les tests sont donc effectués à la volée, pendant que je développe. Comme chaque GameObject peut être exporté sous forme de préfabriqué, il suffit pour tester un objet d'importer le préfabriqué et d'étudier son comportement quand on interagit avec.

    \subsection{Déploiement de l'application}
   
    Le projet n'a pas été déployé sur un serveur réellement, puisqu'il est toujours en développement et que le projet lui-même consiste à faire un prototype de ce que pourrait être un module d'habilitation éléctrique en réalité virtuelle. \\

    Il est à noter qu'un système de build hebdomadaire a été mis en place : mon tuteur souhaite une version utilisable de l'application, toujours à disposition sur le serveur. Chaque fin de semaine, en plus de la sauvegarde du projet sur le serveur, j'effectuais donc un export de l'application au format éxecutable. \\
    
    % Bilan
    \section{Bilan}
    
    \subsection{Bilan personnel}

        Ce stage a été pour moi une véritable source d'enrichissement personnel. Il m'a permit de travailler sur une technologie que j'aime utiliser, dans un cadre de travail très plaisant. \\

        Mes collègues m'ont beaucoup apporté également. Par exemple, un de mes collègues m'a aidé à trouver une alternance pour ma future école d'ingénieur l'année prochaine, en me mettant en relation avec ses contacts. Dans la globalité, le contact avec de nouvelles personnes, des professionnels, a été vraiment un des gros points positifs du stage. \\

        J'ai également eu un aperçu de ce que représente la journée type d'un salarié dans une entreprise. J'ai ainsi compris l'intérêt de la pause du midi, et j'ai eu le confort de pouvoir organiser mes horaires comme je le souhaite, tant que je faisais mes 7 heures quotidiennes. \\

    \subsection{Bilan professionnel}

    Cette experience m'aura également été très bénéfique sur le plan professionnel. En effet, j'ai eu à affronter les difficultés du monde professionnel tout seul, bien que guidé par mon tuteur. J'ai vraiment apprécié cette sensation d'être enfin arrivé au stade où je travaille dans une vraie entreprise, et cela confirme bien mon projet de continuer dans une formation d'ingénieur par apprentissage, afin de rester proche du monde de l'entreprise. \\

    Grâce à ce stage, j'ai également pu apprendre beaucoup de compétences techniques qui vont m'être très utiles dans la suite de mon cursus, notamment dans l'utilisation d'Unity. \\ 

    De plus, j'ai appris énormément sur la communication en entreprise, et les quelques explications techniques que j'ai dû faire auprès de mon tuteur m'ont appris beaucoup sur la vulgarisation qui est nécéssaire pour faire comprendre le développement informatique aux non initiés. \\ 

    J'ai particulièrement apprécié le fait d'avoir un projet concret, et des résultats attendus par des supérieurs hiérarchiques. Ce genre de pression donne envie d'avancer dans le projet. \\
    
    Ce projet concret m'a confronté à la réalité du terrain, et a été un prolongement opérationnel très profitable aux projets d'écoles qui peuvent sembler parfois trop théoriques. Ce stage a donc conclu ma formation de DUT Informatique par une application concrète très appréciable. \\

    % Conclusion
    \huge \textbf{Conclusion} \vspace{5pt} \\
   \normalsize
   
    Ce stage aura été pour moi un vrai succès, autant sur le plan personnel que professionnel. De la montée en compétence à l'intéraction avec mes collèges, la grande majoritée des points sont positifs, et j'ai été très satisfait de mon expérience auprès de Bouygues Energies et Services. \\

    Cette experience m'aura permis de me rendre mieux compte de l'atmosphère de travail dans un grand groupe, et ainsi de valider mon choix d'orientation autour des groupes de plus petite taille, tout du moins dans la suite de ma formation. \\

    La technologie utilisée, le travail que j'ai réalisé, le rythme que j'ai adopté, toutes ces choses ont permis de confirmer mon orientation pour mon projet d'études, mais également mon projet professionnel. Je me suis réellement épanoui dans l'univers où j'ai travaillé, et je souhaite vraiment continuer dans la voie du développement d'applications lourdes, utiliser des moteurs de jeu ou des API graphiques de plus bas niveau. J'ai également découvert mon réel intérêt pour les technologies nouvelles dans l'informatique, comme par exemple la réalité virtuelle. \\

    Le point négatif du stage, la phase de conception qui n'a pas été réalisée correctement, m'aura montré l'intérêt crucial de la définition du cahier des charges préalable à la mise en oeuvre du projet. J'éprouve également quelques regrets dans le fait de ne pas avoir d'application réellement terminée à ajouter à mon CV, et d'avoir fait uniquement un prototype d'une application qui sera certainement terminée par une autre équipe, si les moyens y sont mis. J'aurais vraiment aimé participer à la phase de déploiement et de maintenance du projet, qui pour moi sont des parties tout aussi importantes que la conception et la réalisation. \\

    Bien qu'il reste à l'état de protoype, l'objectif de conception d'un module d'habilitation éléctrique est atteint. Ce projet permettra vraisemblabement d'atteindre une nouvelle étape de production d'une future application qui apportera une véritable plus-value à Bouygues Energies et Services.

    % Annexes
    \section{Annexes}
    
    \newpage 
    % Page de fin 
    \normalsize
    \thispagestyle{empty}
    \noindent
    \begin{minipage}{.5\textwidth}
        Bouygues Energies et Services \\
        49 avenue du Lac du Bourget \\
        73375 Le Bourget du Lac, France
    \end{minipage}
    \begin{minipage}{.5\textwidth}
    \begin{flushright}
        Pinon Olivier \\
        IUT d'Annecy \\
        DUT INFO \\
        Année 2016/2017 \\
    \end{flushright}
    \end{minipage}

    \vspace{10pt}
    \noindent\rule{0.725\paperwidth}{0.4pt}
    
    \vfill 
    \begin{flushleft}
    \huge \textbf{Résumé : } \\
    \vspace{10pt}
    \normalsize Le secteur de l'industrie s'intéresse de plus en plus aux solutions modernes comme la réalité virtuelle pour répondre à leurs besoins. Afin de créer un nouveau module de formation à l'habilitation éléctrique, j'ai travaillé en collaboration avec le bureau d'études de Bouygues Energies et Services, pour développer une application lourde en réalité virtuelle à l'aide du moteur de jeu Unity et du langage C\#. \\
    \end{flushleft}
    
    \vfill 
    \begin{flushleft}
    \huge \textbf{Mots clefs : } \vspace{2pt} \\
    \vspace{10pt}
    \normalsize Bouygues, Module, Habilitation, Electrique, Réalité, Virtuelle, Unity, 3D, C\#
    \end{flushleft}

    \noindent\rule{0.725\paperwidth}{0.4pt}
    
    \vfill 
    \begin{flushleft}
    \huge \textbf{Abstract : } \\
    \vspace{10pt}
    \normalsize There is a growing interest about virtual reality in the industry sector. In order to fulfil the need of a new electrical habilitation training, I worked in collaboration with the technical studies office of Bouygues Energies and Services, to develop a desktop application in virtual reality using the Unity game engine, and the C\# language.
    \end{flushleft}
    
    \vfill 
    \begin{flushleft}
    \huge \textbf{Keywords : } \\
    \vspace{10pt}
    \normalsize Bouygues, Electrical, Habilitation, Training, Virtual, Reality, Unity, 3D, C\#
    \end{flushleft}

	\vfill
    Palanca Jérôme \hfill Damas Luc
\end{document}
